\section{Gaussian initialization--integration identity}
\label{supp:gaussian}
For a centered clean Gaussian coordinate with variance $\lambda>0$, isotropic
variance-preserving noising gives variance $v_t=\bar\alpha_t\lambda+1-\bar\alpha_t$
at timestep $t$, where $\bar\alpha_t$ is the cumulative signal coefficient.
For a value $x_t$ of the noisy coordinate, the exact noise predictor is
$\epsilon^*(x_t,t)=\sqrt{1-\bar\alpha_t}\,x_t/v_t$.
Substituting into deterministic DDIM without clipping gives the multiplier
from timestep $t$ to an earlier timestep $s$:
\begin{equation}
m_{st}=\frac{\sqrt{\bar\alpha_s\bar\alpha_t}\lambda+
\sqrt{(1-\bar\alpha_s)(1-\bar\alpha_t)}}{v_t}.
\label{eq:supp_ddim_step}
\end{equation}
With $u_t=(\sqrt{\bar\alpha_t\lambda},\sqrt{1-\bar\alpha_t})$,
the numerator is $u_s^\top u_t$ and $\|u_t\|^2=v_t$. Thus
$m_{st}=\sqrt{v_s/v_t}\cos\angle(u_s,u_t)$.
Multiplying adjacent grid transitions telescopes the variance factors.
Let $\nu_{\rm start}$ be the actual initial variance, $v_{\rm start}$ and
$v_{\rm end}$ the noising variances at the grid endpoints, and
$X_{\rm out}$ the output coordinate.
\begin{equation}
\frac{\operatorname{Var}(X_{\rm out})}{\lambda}
=\frac{\nu_{\rm start}}{v_{\rm start}}
\frac{v_{\rm end}}{\lambda}
\prod_{t\rightarrow s}\cos^2\angle(u_s,u_t).
\label{eq:supp_ddim_variance}
\end{equation}
The clean endpoint has $v_{\rm end}=\lambda$.
Matched noisy-marginal initialization sets $\nu_{\rm start}=v_{\rm start}$;
standard-normal initialization sets $\nu_{\rm start}=1$.
The numerical oracle uses the checkpoint's actual alpha values and the
scheduler's actual previous-timestep arithmetic. Learned samples keep the
checkpoint's standard clipping behavior, so the exact Gaussian identity is a
surrogate prediction, not an equality asserted for those samples.

For a fixed start and endpoint and ordered angles in the first quadrant,
splitting an angular interval reduces the product-of-cosines deficit because
$\cos a\cos b\geq\cos(a+b)$ for $a,b\geq0$ with $a+b\leq\pi/2$.
This refinement statement does not apply when a changed step count also changes
the initial timestep, when clipping is active, or when the learned dynamics
are nonlinear or noncommuting. Our intermediate-state measurements report
projected off-diagonal mixing and regression residuals to assess these limits.
We do not use a Gaussian integration rule in a covariance-dependent time
coordinate as evidence for a learned-model solver's universal order.
